\section{Intertwining derivatives of matrix powers}
\label{sec:power_derivative_intertwine}

For a positive definite matrix $A$ and $0<p<1$, write
\begin{align*}
  \mathcal D_p(A)[H]
    &=\kappa_p\int_0^\infty s^p(sI+A)^{-1}H(sI+A)^{-1}\,ds.
\end{align*}
Here $\kappa_p$ is the scalar normalization in the integral formula
for the derivative of $A^p$. Set $\mathcal D_0(A)[H]=0$ and
$\mathcal D_1(A)[H]=H$. For Hermitian directions this is the derivative
within the Hermitian matrices; the formula extends complex linearly to
all matrices. This is the derivative formula used in the proof of
Lemma~7.2 of~\cite{OpenAI2026AreaLaw}.

\begin{theorem}[Intertwining the power derivative]
  \label{thm:power_derivative_intertwine}
  \lean{Matrix.rpowDeriv_intertwine}
  \leanok
  Let $A$ and $B$ be positive definite matrices on finite-dimensional
  complex Hilbert spaces, and let $J$ be a linear map satisfying $AJ=JB$.
  For every $p\in[0,1]$ and every matrix $H$ on the space of $B$,
  \begin{align*}
    \mathcal D_p(A)[JHJ^*]&=J\mathcal D_p(B)[H]J^*.
  \end{align*}
\end{theorem}

\begin{proof}
  \leanok
  The endpoints follow from the definitions. For $0<p<1$ and $s>0$,
  both shifted matrices are positive definite. The identity
  $(sI+A)J=J(sI+B)$ gives
  $(sI+A)^{-1}J=J(sI+B)^{-1}$. Taking adjoints gives
  $J^*(sI+A)^{-1}=(sI+B)^{-1}J^*$. Consequently,
  \begin{align*}
    (sI+A)^{-1}JHJ^*(sI+A)^{-1}
      &=J(sI+B)^{-1}H(sI+B)^{-1}J^*.
  \end{align*}
  The linear map $H\mapsto JHJ^*$ is continuous in finite dimension
  and therefore passes through the integrable derivative formula.
  This proves the assertion.
\end{proof}
